%% Đề 45 -- de-tuyen-sinh-lop-10-thpt-mon-toan-nam-2018-2019-so-gddt-lang-son.pdf
\deheader{ĐỀ SỐ 45}%
  {Kỳ thi Tuyển sinh vào lớp 10 THPT — Năm học 2018--2019}%
  {Sở GD\&ĐT Lạng Sơn | Chỉ phần Đại số}
\addcontentsline{toc}{section}{Đề 45 — Tuyển sinh vào lớp 10 THPT — Năm học 2018--2019}

\begin{de_cau}{Câu 1}{2,0 điểm}
\begin{enumerate}[label=\alph*)]
  \item Tính giá trị của biểu thức: $P = \sqrt{81} - \sqrt{25} + \sqrt{9}$.
  \item Rút gọn biểu thức:
  \[Q = \left(\dfrac{1}{\sqrt{a}-1}+\dfrac{1}{\sqrt{a}+1}\right)\cdot\dfrac{a-1}{2}\]
  với $a \ge 0,\; a \neq 1$.
\end{enumerate}
\end{de_cau}

\begin{loi_giai_cau}
\textbf{a)} $P=9-5+3=\boxed{7}$.

\textbf{b)} $\dfrac{1}{\sqrt{a}-1}+\dfrac{1}{\sqrt{a}+1}=\dfrac{(\sqrt{a}+1)+(\sqrt{a}-1)}{a-1}=\dfrac{2\sqrt{a}}{a-1}$.

$Q=\dfrac{2\sqrt{a}}{a-1}\cdot\dfrac{a-1}{2}=\boxed{\sqrt{a}}$.
\end{loi_giai_cau}

\begin{de_cau}{Câu 2}{2,0 điểm}
\begin{enumerate}[label=\alph*)]
  \item Vẽ đồ thị hàm số $y = 2x - 3$ trên mặt phẳng tọa độ $Oxy$.
  \item Giải hệ phương trình:
  \[\begin{cases} 3x-y=5 \\ 2x+3y=7 \end{cases}.\]
\end{enumerate}
\end{de_cau}

\begin{loi_giai_cau}
\textbf{a)} Đồ thị đi qua $(0;-3)$ và $\left(\dfrac{3}{2};0\right)$. Vẽ đường thẳng qua hai điểm này.

\textbf{b)} Nhân pt(1) với 3: $9x-3y=15$. Cộng pt(2): $11x=22\Rightarrow x=2$. Từ (1): $y=1$.

Nghiệm $\boxed{(x;y)=(2;1)}$.
\end{loi_giai_cau}

\begin{de_cau}{Câu 3}{2,0 điểm}
Cho phương trình: $x^2 - 2mx + 2m - 1 = 0$ (với $m$ là tham số).
\begin{enumerate}[label=\alph*)]
  \item Giải phương trình khi $m = 2$.
  \item Chứng minh rằng phương trình luôn có nghiệm với mọi giá trị của $m$.
  \item Tìm $m$ để phương trình có hai nghiệm $x_1, x_2$ thỏa mãn $x_1^2 + x_2^2 = 10$.
\end{enumerate}
\end{de_cau}

\begin{loi_giai_cau}
\textbf{a)} $m=2$: $x^2-4x+3=0\Rightarrow(x-1)(x-3)=0\Rightarrow\boxed{x_1=1,\;x_2=3}$.

\textbf{b)} $\Delta'=m^2-(2m-1)=(m-1)^2\ge0$ với mọi $m$. Phương trình luôn có nghiệm. $\square$

\textbf{c)} Vi-ét: $x_1+x_2=2m$, $x_1x_2=2m-1$.

$x_1^2+x_2^2=(2m)^2-2(2m-1)=4m^2-4m+2=10\Rightarrow4m^2-4m-8=0\Rightarrow m^2-m-2=0\Rightarrow\boxed{m=2}$ hoặc $\boxed{m=-1}$.
\end{loi_giai_cau}

\begin{de_cau}{Câu 5}{0,5 điểm}
Cho các số thực không âm $x, y$ thỏa mãn $x + y = 1$. Tìm giá trị lớn nhất và giá trị nhỏ nhất của biểu thức:
\[S = x^3 + y^3 + 2xy.\]
\end{de_cau}

\begin{loi_giai_cau}
$S=(x+y)(x^2-xy+y^2)+2xy=(x^2+y^2-xy)+2xy=x^2+y^2+xy=(x+y)^2-xy=1-xy$.

Vì $x,y\ge0$ và $x+y=1$: $0\le xy\le\dfrac{1}{4}$.

$S=1-xy\in\left[\dfrac{3}{4};1\right]$.

$S_{\min}=\boxed{\dfrac{3}{4}}$ khi $x=y=\dfrac{1}{2}$; $S_{\max}=\boxed{1}$ khi $(x;y)=(1;0)$ hoặc $(0;1)$.
\end{loi_giai_cau}
